\documentclass[12pt]{article}
\usepackage[utf8]{inputenc}
\usepackage[T1]{fontenc}
\usepackage{hyperref}
\title{Sula Language Resources for K-12 Curriculum Development}
\author{Compiled by Claude Sonnet 5.5 (claude-sonnet-5-5) for Andr\'as Kornai}
\date{}
\begin{document}
\maketitle

\section{Introduction}
Sula (ISO 639-3 szn) is an Austronesian language of the Central Maluku branch, related to Buru. It is spoken mainly on Sanana island and in the Sula archipelago of North Maluku, Indonesia. Mangole (mqc) has its own code and is sometimes treated as a Sula dialect. Speaker estimates conflict: Wikipedia cites 20,000 native speakers (1983), and a 2020 University of Hawai\textquotesingle i dissertation says up to 47,000 residents, describing the language as critically to severely endangered. Sources disagree on dialects (three or four, or perhaps two main ones). There is heavy contact influence from Ambonese Malay, Ternate and Indonesian.

\section{Primary Educational Resources}
None found.

\section{Linguistic and Reference Materials}
\begin{itemize}
\item \emph{Sula: Its language, land, and people}, a 2020 University of Hawai\textquotesingle i dissertation: \url{https://scholarspace.manoa.hawaii.edu/items/daeb65e2-8b98-4002-93b2-ced62f91fe4d}
\item Wikipedia article on the Sula language: \url{https://en.wikipedia.org/wiki/Sula_language}
\end{itemize}

\section{Religious and Community Resources}
\begin{itemize}
\item Joshua Project profile of the Sula in Indonesia: \url{https://joshuaproject.net/people_groups/14718/ID}
\item Asia Harvest people-group profile: \url{https://www.people-groups.asiaharvest.org/Indonesia/Sula.pdf}
\end{itemize}

\section{Digital and Audio Resources}
None found in these searches.

\section{Curriculum Development Context}
The language is not known to be taught in schools. The only substantial resource found is the scholarly dissertation, which a revitalization programme could use as a documentary base. Pressure from Malay and Indonesian is strong. Local community and education bodies were not found and should be identified through further work.
\end{document}